\subsection{Orthogonalität und Winkel}
\label{subsec:vector_algebra_orthogonality_angles}

\begin{KeyPointBox}{Orthogonalitätskriterium}{vector_algebra_orthogonality_criterion}
  Zwei von null verschiedene Vektoren sind genau dann
  \emph{orthogonal} \EnglishTerm{orthogonal}, wenn ihr Skalarprodukt null ist:
  \[
    \boxed{\vec a\perp\vec b\quad\Longleftrightarrow\quad
      \vec a\cdot\vec b=0}.
  \]
\end{KeyPointBox}

\begin{ExampleBox}{Eine orthogonale Richtung bestimmen}{vector_algebra_find_orthogonal_vector}
  Für
  \[
    \vec a=\begin{pmatrix}1\\2\end{pmatrix},
    \qquad
    \vec b=\begin{pmatrix}4\\x\end{pmatrix}
  \]
  soll $\vec a\perp\vec b$ gelten. Dann
  \[
    \vec a\cdot\vec b=0
    \quad\Longleftrightarrow\quad
    4+2x=0
    \quad\Longleftrightarrow\quad
    x=\Cloze{2cm}{-2}.
  \]
  \begin{center}
    \CachedTikzFigure{content/resources/chapters/vector_algebra/tikz/orthogonal_vectors_example.tex}
  \end{center}
\end{ExampleBox}

\begin{KeyPointBox}{Winkel zwischen zwei Vektoren}{vector_algebra_angle_between_vectors}
  Für $\vec a,\vec b\ne\vec 0$ folgt aus der geometrischen Definition
  \[
    \cos(\varphi)=\frac{\vec a\cdot\vec b}{|\vec a|\,|\vec b|}
    \qquad\Longrightarrow\qquad
    \boxed{
      \varphi=\arccos\!\left(
        \frac{\vec a\cdot\vec b}{|\vec a|\,|\vec b|}
      \right)
    }.
  \]
\end{KeyPointBox}

\begin{ExampleBox}{Winkel zweier Vektoren}{vector_algebra_angle_example}
  \noindent
  \begin{minipage}[c]{0.56\linewidth}
    Für $\vec a=(1,2)^{\mathsf T}$ und
    $\vec b=(4,0)^{\mathsf T}$ gilt
    \[
      \vec a\cdot\vec b=4,
      \quad |\vec a|=\sqrt5,
      \quad |\vec b|=4,
    \]
    \[
      \varphi
      =\arccos\!\left(\frac{4}{4\sqrt5}\right)
      =\Cloze{4cm}{\arccos\!\left(\frac1{\sqrt5}\right)
        \approx63{,}4^\circ}.
    \]
  \end{minipage}\hfill
  \begin{minipage}[c]{0.40\linewidth}
    \centering
    \CachedTikzFigure{content/resources/chapters/vector_algebra/tikz/angle_between_vectors_example.tex}
  \end{minipage}
\end{ExampleBox}
